% Generated by tools/edn_latex.py from reports/edn.md.
% Do not edit: change reports/edn.md and run `make edn`.
\ednentry{
  id = {EDN-54},
  title = {Tracking degrades silently when it is not configured, with an environment variable to make it fatal},
  date = {2026-09-30},
  milestone = {M3: Quality Assurance},
  topic = {CI/CD, Experiment Tracking},
  participants = {@lukas2510, decided by the repository owner},
  decision = {A pipeline stage opens its MLflow run through \texttt{optional\_\allowbreak{}run}, which yields a real run handle when the server is configured and reachable and a \texttt{DisabledRun} that logs nothing when it is not -- whether because nothing is configured, because the credentials are refused or because the server cannot be reached. The stage runs to completion either way and records \texttt{mlflow.\allowbreak{}tracking\_\allowbreak{}mode} in \texttt{model.\allowbreak{}json}. Setting \texttt{RECOMMENDITOS\_\allowbreak{}REQUIRE\_\allowbreak{}TRACKING=\allowbreak{}1} turns any of those failures back into a failed stage. There is never a fallback to a local tracking store.},
  rationale = {\emph{Alternatives considered:}
\begin{itemize}[leftmargin=*, topsep=2pt]
\item \textbf{Option A: fail the stage whenever tracking is not available,} which is what \texttt{tracked\_\allowbreak{}run} does and what EDN-46's helper was built for.
\newline Pros: a run that was not recorded never happened, which is the strongest possible reading of NFR-06; and the failure is the cheapest one, before the fit.
\newline Cons: CI has no DagsHub credentials and a fresh clone has none either, so \texttt{pytest} and \texttt{dvc repro} would both fail on the one path every contributor and every pull request takes. The suite's whole premise, stated in \texttt{tests/\allowbreak{}conftest.\allowbreak{}py}, is that it runs on a clean clone with no credentials and no network; this option contradicts it.
\item \textbf{Option B: fall back to a local tracking store when the server is not configured.}
\newline Pros: the run is recorded somewhere, so nothing is lost and the code path is the same.
\newline Cons: since MLflow 3.16 an unset tracking URI resolves to \texttt{sqlite:/\allowbreak{}/\allowbreak{}/\allowbreak{}\$PWD/\allowbreak{}mlflow.\allowbreak{}db}, and for a DVC stage \texttt{\$PWD} is the repository root. The fallback therefore drops a database into the repository on every run of every contributor, gitignored so nothing would ever point it out. \texttt{tests/\allowbreak{}conftest.\allowbreak{}py} already guards \texttt{mlflow.\allowbreak{}db} and \texttt{mlruns/\allowbreak{}} as forbidden paths for exactly this reason, and writing the guard for \texttt{models/\allowbreak{}} caught a real instance of it: a local store puts its \emph{artifacts} in \texttt{.\allowbreak{}/\allowbreak{}mlruns} too.
\item \textbf{Option C (chosen): degrade silently, with \texttt{RECOMMENDITOS\_\allowbreak{}REQUIRE\_\allowbreak{}TRACKING=\allowbreak{}1} as the escape hatch.}
\newline Pros: every contributor and CI can train, which is a property of the whole pipeline rather than of this stage. And the runs whose numbers are cited can be made to fail loudly, so ``the runs are on the server'' is something a person can verify instead of something a reader has to take on trust.
\newline Cons: two behaviours instead of one, and the quiet one is the default. A contributor who \emph{meant} to record a run and mistyped a variable gets a warning line in a long stage log and a model that looks finished. The mitigation is that \texttt{model.\allowbreak{}json} carries \texttt{tracking\_\allowbreak{}mode: \textquotedbl{}disabled\textquotedbl{}} and a null \texttt{run\_\allowbreak{}id}, so the artefact itself says it was not recorded, and \texttt{evaluate} says so again when it finds no run to resume.
\item \textbf{Option D: degrade silently with no escape hatch.}
\newline Pros: one behaviour, nothing to remember.
\newline Cons: it makes the report's claim unfalsifiable. Nobody could demonstrate that a given set of runs reached the server, because the successful and the skipped path are indistinguishable from the outside.
\end{itemize}
\par
\emph{Why:} The two halves of the decision answer two different audiences and neither can be dropped.
\par
Silence is right for the default audience, which is CI and a contributor on a fresh clone. Training with no credentials is not a degraded mode for this project; it is the normal mode for everything except the handful of runs that produce numbers for the report, and a gate that turns red for everyone stops being read.
\par
The escape hatch is right for the other audience, which is whoever produces those numbers. It was not hypothetical: the live verification of this stage against the team's DagsHub server was run with \texttt{RECOMMENDITOS\_\allowbreak{}REQUIRE\_\allowbreak{}TRACKING=\allowbreak{}1} precisely so that a misconfiguration could not let the check pass while logging nothing. Without it the verification would have proved that the stage finishes, which it does in either case, rather than that the runs landed.
\par
The rule that there is never a local fallback is what makes ``no tracking URI means tracking is off'' a statement with content, and it is the same reasoning \texttt{recommenditos/\allowbreak{}config.\allowbreak{}py} already applies to loading \texttt{.\allowbreak{}env} from a named path rather than searching upwards.
\par
One consequence was measured rather than assumed, and it is recorded in EDN-57: detecting ``the server is not reachable'' is not free, because MLflow retries. That is why the retry budget is pinned, and the two decisions have to be read together.},
  ai = {Alternative generation, Alternative assessment, Recommendation, Solution generation},
  response = {Accepted},
  assessment = {AI identified that the existing helper's fail-fast behaviour, which EDN-46 chose deliberately for a credential check, is the wrong behaviour for a pipeline stage, and that the obvious repair -- a local fallback -- would put a database in the repository root under MLflow 3.16's changed default. It proposed the null-object handle so the branch exists once rather than in front of every log call, and the environment variable so the quiet default stays falsifiable. It also used the variable in its own live verification rather than only testing it, which is what showed the hatch earns its place.},
  aievidence = {Claude Code session on 2026-09-30 implementing issue \#37: the DagsHub verification was run with \texttt{RECOMMENDITOS\_\allowbreak{}REQUIRE\_\allowbreak{}TRACKING=\allowbreak{}1} and the throwaway experiment deleted afterwards; the \texttt{mlruns} leak was found by the new \texttt{models/\allowbreak{}} guard in \texttt{tests/\allowbreak{}conftest.\allowbreak{}py} failing.},
  evidence = {\href{https://github.com/mlops-2627q1-mds-upc/MLOPS_Recommenditos/blob/main/recommenditos/tracking.py}{\texttt{recommenditos/\allowbreak{}tracking.\allowbreak{}py}}, \texttt{optional\_\allowbreak{}run}, \texttt{Run} and \texttt{DisabledRun}; \texttt{tests/\allowbreak{}test\_\allowbreak{}model.\allowbreak{}py::test\_\allowbreak{}a\_\allowbreak{}tracking\_\allowbreak{}failure\_\allowbreak{}does\_\allowbreak{}not\_\allowbreak{}fail\_\allowbreak{}training}, \texttt{::test\_\allowbreak{}require\_\allowbreak{}tracking\_\allowbreak{}makes\_\allowbreak{}a\_\allowbreak{}failure\_\allowbreak{}fatal}, \texttt{::test\_\allowbreak{}tracking\_\allowbreak{}is\_\allowbreak{}disabled\_\allowbreak{}without\_\allowbreak{}a\_\allowbreak{}tracking\_\allowbreak{}uri} and \texttt{::test\_\allowbreak{}require\_\allowbreak{}tracking\_\allowbreak{}refuses\_\allowbreak{}a\_\allowbreak{}model\_\allowbreak{}with\_\allowbreak{}no\_\allowbreak{}run\_\allowbreak{}to\_\allowbreak{}resume}; the \texttt{FORBIDDEN\_\allowbreak{}PATHS} guard in \href{https://github.com/mlops-2627q1-mds-upc/MLOPS_Recommenditos/blob/main/tests/conftest.py}{\texttt{tests/\allowbreak{}conftest.\allowbreak{}py}}; \hyperref[EDN-46]{EDN-46}; \hyperref[EDN-57]{EDN-57}; \href{https://github.com/mlops-2627q1-mds-upc/MLOPS_Recommenditos/issues/37}{issue \#37}.},
}
